% Shared module: recovering the components u_j, interference, recursion, labels, acceptance.
\subsection{Recovering the groups as functions}\label{sec:alg-components}
Once modes are assigned, the component of group $j$ is the sum of its modes,
\begin{equation}\label{eq:uj}
  u_j(t,x)=\sum_{p:\,j(p)=j}A_p\,e^{\lambda_p(t-t_0)+i\kappa_px},
\end{equation}
the discrete counterpart of the projector~\eqref{eq:lagrange}. A mode within distance $3$ of two
varieties lies in the overlap counted by $\deltaP$ (for the running example, the $\kappa=0$ mode); it is
split equally between the groups. This is the minimum-norm, or minimum-energy, choice among all splits allowed by
\cref{thm:factor}: if the shared mode has amplitude $A$ and is given as $\alpha A$ to one group and $(1-\alpha)A$ to the
other, the energy it adds to the components is $\abs{A}^2\big(\alpha^2+(1-\alpha)^2\big)$, which is smallest at
$\alpha=\tfrac12$. The data cannot tell the splits apart, so the algorithm takes the one that puts the least energy
into the ambiguity. It coincides with the posterior mean of the Bayesian layer (\cref{sec:bayes-split}), so the
convention is not arbitrary. The energy shares $w_j$ of~\eqref{eq:energy}
are computed from these components.

\subsection{Interference between groups and the recursion over components}\label{sec:alg-interference}
The groups are not orthogonal. Let $B_j$ be an orthonormal basis of the sampled kernel of $Q_j$ (all
exponentials $e^{\lambda t+i\kappa x}$ with $Q_j(\lambda,i\kappa)=0$, $\kappa\in\mathcal K$, on the grid) and
$\sumop=[B_1\cdots B_N]$. Then $\gN$ is the smallest eigenvalue of $\sumop^H\sumop$ above $10^{-8}$, the
eigenvalues below that threshold count the directions shared by two groups, and for two groups
$\gamma_2=1-\cF$. Because pairwise angles do not determine $\gN$ for $N\ge3$ (\cref{sec:cond}), the full
Gram matrix is used.

Each component $u_j$ is then analysed again by the same procedure, up to depth two: a factor that still
splits on its own component is replaced by its children. A split is accepted without further evidence
if it leaves the kernel unchanged, and otherwise only if it improves the evidence of \cref{sec:alg-accept} by more
than one nat. Complex factors of a real field are always kept together with their conjugates. The beam equation of
\cref{app:cases} is an example. The core first fits it as one factor of degree two in $\tau$,
$\tau^2+0.0903\,\xi^4$, whose component $u_1$ contains, at every wavenumber, two modes with rates
$\lambda=\pm0.30\,i\kappa^2$. The recursion finds that this factor splits along its two root branches, into the children
$\tau-0.3004\,i\xi^2$ and $\tau+0.3004\,i\xi^2$; their kernels together are the kernel of the parent, so the split is
accepted without further evidence, and the two children, which are complex conjugates, are reported together.

\subsection{Identifiability and the output labels}\label{sec:alg-labels}
\Cref{sec:certify} rules out the label ``identified''. Each reported structure carries one of four labels.
\begin{description}[leftmargin=2.9cm,labelwidth=2.7cm]
  \item[\FC] Every factor $Q_j$, of degree $m_j$ in $\tau$ and $d_j$ in $\xi$, has at least $m_j(d_j+1)+1$ modes
  (more data than coefficients), modes at $d_j+1$ distinct wavenumbers (\cref{prop:counting}), and $m_j$ modes
  at some single wavenumber (one per root); and the fit passes the goodness-of-fit test below.
  The data are compatible with the factored equation and pin down its coefficients within the class.
  \item[\IDK] The fit is good, but some factor is determined only relative to the declared class (too
  few distinct wavenumbers): another equation of a larger class fits equally well.
  \item[\MIS] The goodness-of-fit test rejects the structure.
  \item[\UND] The data cannot decide, typically because of confluent rates or a repeated factor.
\end{description}
Repeated factors ($Q=P_i^2$) are always reported as \UND: by \cref{thm:ident}(b) a repeated factor is
visible only through polynomial prefactors $t\,e^{\lambda t}$, which finite noisy data cannot separate from
two nearby rates.

\subsection{Noise level and acceptance}\label{sec:alg-accept}
The noise level enters in three places. The observation noise $\noise$ is integrated out in the
interpolation (\cref{sec:alg-interp}). Per Fourier bin it becomes $\sigbin$, which sets the band and the
rate uncertainties. In the final decision it is used against the \emph{raw observations}, not against the
interpolated field. A candidate set of factors $Q_1,\dots,Q_N$ defines a linear model for the vector $y$ of the $n$
observations: let $E_j$ be the matrix whose columns are the exponentials $e^{\lambda t+i\kappa x}$ of the kernel of $Q_j$
(the functions spanned by $B_j$ in \cref{sec:alg-interference}), evaluated at the observation points, $E=[E_1\cdots E_N]$,
and let $a$ be the vector of their amplitudes. Then
\[
  y=Ea+\varepsilon,\qquad \varepsilon\sim\Normal(0,\noise^2I),\qquad a\sim\Normal(0,\omega^2I),
\]
where $I$ is the identity matrix, so that $\noise^2I$ and $\omega^2I$ are the covariance matrices of the noise and of the
amplitudes, and $\omega^2$ is the prior variance of each amplitude. Both are multivariate Gaussian distributions, and
integrating the amplitudes out gives the distribution of the data under the candidate,
\[
  y\mid Q_1,\dots,Q_N\ \sim\ \Normal\big(0,\ \noise^2I+\omega^2EE^H\big),
\]
with $E^H$ the conjugate transpose of $E$. The score of the candidate adds a complexity penalty to the logarithm of
this density,
\begin{equation}\label{eq:raw-evidence}
  S(Q_1,\dots,Q_N)=\log\Normal\big(y;\,0,\ \noise^2I+\omega^2EE^H\big)-\tfrac12\,n_{\rm par}\log n ,
\end{equation}
where $n_{\rm par}$ is the number of coefficients that define the factors. The first term is the log-probability of
the observations actually made under the candidate factors, with the amplitudes integrated out; the second, as in the
Bayesian information criterion, approximately accounts for the coefficients of the factors themselves, which were fitted
rather than integrated, and penalizes candidates with more of them.
The score is built to \emph{compare} candidate structures on the same data, which of two factor sets explains the
observations better, not to judge one fit in isolation. It decides between candidates: dropping a factor, dropping a term whose $90\%$
interval contains zero, merging two factors, or splitting one. A change is accepted only if it gains more
than one nat. The same model gives a goodness-of-fit test: with $\hat y$ the projection of $y$ on the columns of $E$,
\[
  \chi^2=\sum_{i=1}^{n}\frac{\abs{y_i-\hat y_i}^2}{\noise^2},
\]
which, if the candidate is right and the noise Gaussian (variance $\noise^2$ per real observation), approximately follows
a chi-squared distribution over noise realizations with $n-\rank E$ degrees of freedom; $\rank E$ rather than the number of
columns, because a mode shared by two groups appears in both $E_1$ and $E_2$, and only approximately, because the rates in
$E$ were themselves fitted to the data. (For complex data with $\mathbb E\abs{\varepsilon_i}^2=\noise^2$ the statistic is half
a chi-squared variable with $2(n-\rank E)$ degrees of freedom.) A $p$-value below $10^{-3}$
turns the label into \MIS. Deciding on the raw data matters:
the interpolant has already smoothed the high wavenumbers, and a candidate that fits the interpolant may
not fit the observations.

Finally the interaction ratio $\rho$ of~\eqref{eq:energy} is computed for every evolution factor (degree one in $\tau$) with more than one
term from the closed form~\eqref{eq:eint-fourier}. It is reported, not used to accept or reject: by
\cref{prop:discont} it is a diagnostic of the evolution inside a factor, not a test of the factor.
